%----------------------------------------------------------------------------------------------
%--------------------------------     3. Kurvendiskussionen    --------------------------------
%----------------------------------------------------------------------------------------------
\section{Kurvendiskussionen}


%----------------------------------------------------------------------------------------------
%---------------------------------     3.1 Satz von Rolle    ----------------------------------
%----------------------------------------------------------------------------------------------
\subsection{Satz von Rolle}{454}

\begin{tabular}{ll}
    Voraussetzungen:    & $f$ auf Intervall [$a;b$] mind. $(n+1)$ mal ableitbar \\
                        & $f(a) = f(b)$
\end{tabular}
\vspace{0.2cm}

Auf dem Intervall ($a;b$) existiert \textbf{mindestens einmal} eine horizontale Tangente: $f'(\xi) = 0$		

\begin{center}
    \includegraphics[width=0.9\columnwidth]{images/rolle.PNG}
\end{center}


%----------------------------------------------------------------------------------------------
%---------------------------------     3.2 Mittelwertsatz    ----------------------------------
%----------------------------------------------------------------------------------------------
\subsection{Mittelwertsatz}{454}

\begin{tabular}{ll}
    Voraussetzung:    & $f$ auf Intervall [$a;b$] mind. $(n+1)$ mal ableitbar   
\end{tabular}

\begin{minipage}[c]{0.5\columnwidth} 
    \includegraphics[width=\columnwidth]{images/mittelwertsatz.PNG}
\end{minipage}
\hfill
\begin{minipage}[c]{0.45\columnwidth} 
\textbf{Sekantensteigung:}\\
$m_s = \frac{\Delta y}{\Delta x} = \frac{f(b) - f(a)}{b - a}$ 

\vspace{0.5cm}

\textbf{Tangentensteigung /} \\
\textbf{Durchschnittssteigung:}\\
$f'(\xi) = \frac{f(b) - f(a)}{b - a} \Rightarrow t_\xi \parallel s$\\
	
\end{minipage}			


%----------------------------------------------------------------------------------------------
%------------------------------     3.3 Monotoniebedingungen    -------------------------------
%----------------------------------------------------------------------------------------------
\subsection{Monotoniebedingungen}{453}

Ist $ f(x) $ auf einem zusammenhängenden Intervall stetig und im Inneren differenzierbar, so gilt:\\
\renewcommand{\arraystretch}{1.25}
\begin{tabular}{ll}
    $\Rightarrow f'(x_0) \geq 0$    & für eine monoton wachsende Funktion \\
    $\Rightarrow f'(x_0) > 0$    & für eine streng monoton wachsende Funktion \\
    $\Rightarrow f'(x_0) \leq 0$    & für eine monoton fallende Funktion \\
    $\Rightarrow f'(x_0) < 0$    & für eine streng monoton fallende Funktion 
\end{tabular}


%----------------------------------------------------------------------------------------------
%---------------------------------     3.4 Extremastellen    ----------------------------------
%----------------------------------------------------------------------------------------------
\subsection{Extremastellen}{455}

\subsubsection{Relative Extremalstelle}

\renewcommand{\arraystretch}{1.2}
\begin{tabular}{llll}
    Lokales Maximum:    & & &'Berg' (nicht am Rand der Funktion)               \\ 
                        & & &Wechsel von $\uparrow$ zu $\downarrow$ \\
    \hline
    Lokales Minimum:    & & &'Tal' (nicht am Rand der Funktion)                \\
                        & & & Wechsel von $\downarrow$ zu $\uparrow$ \\
    \hline
    Terrasse:           & & &$f'(x_0) = 0$ aber kein 'Berg' / 'Tal' 
\end{tabular}\\

\renewcommand{\arraystretch}{1}
\textbf{Extremstellen}\\
\begin{tabular}{llll}
    1. &$f'(x)$ bilden \\
    2. &$f'(x)=0$ lösen & $\quad \quad  \Longrightarrow \quad$ &   $\;\;\,$Kandidaten $x_0$\\
    3. &$f''(x_0)$ berechnen & $\quad \quad  \Longrightarrow \quad$ &   
    \begin{tabular}{l}
        $f''(x_0) > 0 \Rightarrow \text{Minimum}$\\
        $f''(x_0) < 0 \Rightarrow \text{Maximum}$\\
        $f''(x_0) = 0 \Rightarrow \text{keine Aussage}$
    \end{tabular}
\end{tabular}

%----------------------------------------------------------------------------------------------
%----------------------------------     3.5 Wendestelle    ------------------------------------
%----------------------------------------------------------------------------------------------
\subsection{Wendepunkt(Terrassenpunkte)}{457}

Wendepunkt = Krümmungswechsel = Vorzeichenwechsel für $f''(x)$ bei $x_0$ 
\vspace{0.2cm}

\textbf{Extremstellen}\\
\begin{tabular}{llll}
    1. &$f''(x)$ bilden \\
    2. &$f''(x)=0$ lösen & $\quad \quad  \Longrightarrow \quad$ &   $\;\;\,$Kandidaten $x_0$\\
    3. &$f'''(x_0)$ berechnen & $\quad \quad  \Longrightarrow \quad$ &   
    \begin{tabular}{l}
        $f'''(x_0) > 0 \Rightarrow$ Wechsel von $\uparrow$ zu $\downarrow \quad$  \textbackslash \\
        $f'''(x_0) < 0 \Rightarrow$ Wechsel von $\downarrow$ zu $\uparrow\quad$ / \\
        $f'''(x_0) = 0 \Rightarrow \text{keine Aussage}$
    \end{tabular}
\end{tabular}


%----------------------------------------------------------------------------------------------
%----------------------------------     3.6 bsp    ------------------------------------
%----------------------------------------------------------------------------------------------
\subsection{Bsp: Extrema und Wendepunkt}

Gegeben ist die Funktion\\
\[f(x) = x^3-4x^2+3x+1\]
Untersuchen Sie die Funktion vollständig hinsichtlich ihrer besonderen Punkte.\\
\begin{tabular}{ll}
    1. & Bestimmen Sie die Extremstellen und geben Sie an, ob es sich um ein \\
       & Maximum oder Minimum handelt.\\
    2. & Bestimmen Sie die Wendestelle(n) der Funktion.
\end{tabular}

\vspace{0.3cm}
\renewcommand{\arraystretch}{2}
\begin{tabular}{lllll}
    $f(x) = x^3-4x^2+3x+1$ & $\Rightarrow$ & $f'(x) = 3x^2-8x+3$ & $\Rightarrow$ & $f''(x) = 6x-8$
\end{tabular} 
\begin{tabular}{lllll}
    $\Rightarrow$ & $f'''(x) = 6$
\end{tabular} 
\vspace{2.5mm}
\hrule
\vspace{2mm}

\begin{tabular}{lllll}
    $f'(x) = 0$ &  $\Rightarrow$ & $3x^2-8x+3 = 0$ & $\Rightarrow$ & 
    \begin{tabular}{lll}
        $\frac{8+\sqrt{28}}{6}$ & $\Rightarrow$ & $x_1 = \frac{4+\sqrt{7}}{3}$\\
        $\frac{8-\sqrt{28}}{6}$ & $\Rightarrow$ & $x_2 = \frac{4-\sqrt{7}}{3}$
    \end{tabular} 
\end{tabular} 

\begin{tabular}{lllllll}
    $f''(x_1) = 6x_1-8$ &  $\Rightarrow$ & $6 \cdot \frac{4+\sqrt{7}}{3}-8$ &  $\Rightarrow$ & $2\sqrt{7} > 0$ & $\Rightarrow$ & $x_1 = x_{min}$ \\
    $f''(x_2) = 6x_2-8$ &  $\Rightarrow$ & $6 \cdot \frac{4-\sqrt{7}}{3}-8$ &  $\Rightarrow$ & $-2\sqrt{7} < 0$ & $\Rightarrow$ & $x_2 = x_{max}$ 
\end{tabular}
\vspace{2.5mm}
\hrule
\vspace{2mm}

\begin{tabular}{lllllll}
    $f''(x) = 0$   &  $\Rightarrow$ & $6x-8 = 0$ &  $\Rightarrow$ & $x_0 = \frac{4}{3} \Rightarrow$ Wendepunkt  \\
    $f'''(x) = 6$  &  $\Rightarrow$ & $6 > 0$ &  $\Rightarrow$ &  von $x_{max}$ zu $x_{min}$
\end{tabular} 



%----------------------------------------------------------------------------------------------
%-------------------------------     3.7 Bernoulli-Hôpital    ---------------------------------
%----------------------------------------------------------------------------------------------
\subsection{Bernoulli-Hôpital}{57-58}

\subsubsection{B.H. I}

$$ \lim_{x \to x_0} \frac{f_1(x)}{f_2(x)} = \frac{0}{0} $$
Wenn $\frac{f'_1(x)}{f'_2(x)}$ eine bestimmte Form ist, dann ist dies das Resultat von 
$\lim \limits_{x \to x_0} \frac{f_1(x)}{f_2(x)}$


\subsubsection{B.H. II}	
$$ \lim_{x \to x_0} \frac{f_1(x)}{f_2(x)} = \frac{\infty}{\infty}$$
Wenn $\frac{f'_1(x)}{f'_2(x)}$ eine bestimmte Form ist, dann ist dies das Resultat von 
$\lim \limits_{x \to x_0} \frac{f_1(x)}{f_2(x)}$

\vspace{0.5cm}
\textbf{Bernoulli-Hôpital darf auch mehrfach nacheinander verwendet werden}
\textbf{\textrightarrow\ immer erst algebraisch verienfachen!} 


%----------------------------------------------------------------------------------------------
%--------------------------------     3.8 Bernoulli-Formen    ---------------------------------
%----------------------------------------------------------------------------------------------
\subsection{Unbestimmte Formen zu Bernoulli-Formen}		

\renewcommand{\arraystretch}{1.5}
\begin{tabular}{ll}
    $f \cdot g$ vom Typ $(0^+) \cdot \infty$        & $\frac{f}{1 / g}$ von Typ $\frac{0}{0}$ \\
                                                    & $\frac{1 / f}{g}$ von Typ $\frac{\infty}{\infty}$ \\
    \\
    $f - g$ von Typ $\infty - \infty$               & $\frac{1/g - 1/f}{1/ fg}$ von Typ $\frac{0}{0}$   \\
    \\
    $f^g$ als $(0+)^0$; $\infty^0$; $1^{\infty}$    & $f^g = e^{g \cdot \ln(f)}$ wobei $g \cdot \ln(f)$ von Typ \\
                                                    & $(0+) \cdot \infty$ bzw. $ \infty \cdot 0$ \\
    \\
    Vorzeichen ausklammern:                         & $(0-) \cdot \infty = -(0+) \cdot \infty$ \\
                                                    & $(0+)^{0-} = \frac{1}{(0+)^{0+}}$ oder $1^{-\infty} = \frac{1}{1^{\infty}}$
\end{tabular}


%----------------------------------------------------------------------------------------------
%---------------------------------     3.9 Polynomdivision     --------------------------------
%----------------------------------------------------------------------------------------------
\subsection{Polynomdivision }{15}

Liefert Summe aus \textbf{ganzrationaler Funktion} und \textbf{echt gebrochener Funktion}.
\[
R(x) = \frac{P_4(x)}{Q_2(x)} = \frac{3x^4 - 10x^3 + 22x^2 - 24x + 10}{x^2 - 2x - 3}, \quad n > m \;\; \Rightarrow \text{unecht gebrochen}
\]

\[
\makebox[270pt][l]{%
\begin{tabular}{@{}l@{}}
$(3x^4 - 10x^3 + 22x^2 - 24x + 10) \; : \; (x^2 - 2x + 3)$ \\[1pt]

$\;3x^2 \cdot (x^2 - 2x + 3)$ \\[-6pt]
\rule{3.5cm}{0.4pt} \\[-2pt]

\hspace*{0.6cm}$-\;4x^3 + 13x^2 - 24x + 10$ \\[1pt]

\hspace*{0.6cm}$-\; 4x \cdot (x^2 - 2x + 3)$ \\[-6pt]
\quad \quad \,\rule{2.9cm}{0.4pt} \\[-2pt]

\hspace*{1.7cm}$5x^2 - 12x + 10$ \\[1pt]

\hspace*{1.7cm}$5 \cdot (x^2 - 2x + 3)$ \\[-6pt]
\quad \quad \quad \quad \quad \,\rule{2.1cm}{0.4pt} \\[-2pt]

\hspace*{2.45cm}$-\;2x - 5 \quad \quad \quad\Rightarrow \quad \quad \quad R(x) = 3x^2 - 4x + 5 + \frac{2x - 5}{x^2 - 2x + 3}$ \\[1pt]
\end{tabular}}
\]

%----------------------------------------------------------------------------------------------
%--------------------------------     3.10 Asymptotenbestimmung     ---------------------------
%----------------------------------------------------------------------------------------------
\subsection{Asymptoten bestimmung}

\label{Asymptotenbestimmung}

Asymptote einer gebrochen rationalen Funktion $r(x) = \frac{P_n(x)}{Q_m(x)}$ bestimmen gem"ass: 

\renewcommand{\arraystretch}{1.3}
\begin{tabular}{l | c | c | c }
	& $m < n$ 								& $m = n$ 	& $n < m$ \\
	\hline
	$\lim\limits_{x \to \pm \infty} r(x)$ 	& $\infty$ oder $-\infty$ 		& $\frac{a_n}{b_m}$ 				& 0 \\
	Asymptote 								& ganzrat. Teil der 	        & Parallel zur $x$-Achse 		    & x-Achse \\
											& Polynomdivision			    & $y = g(x) = \frac{a_m}{b_n}$ 		& 
\end{tabular}
\renewcommand{\arraystretch}{1}
\vspace{0.4em}

\subsubsection*{Bsp. Asymptotenbestimmung}

\begin{tabular}{l}
    $ \displaystyle{ m < n:   \frac{x^2}{x-1} \quad \Rightarrow \quad x + \frac{x}{x-1} \quad \Rightarrow \quad x + \frac{x-1}{x-1} + \frac{1}{x-1} = x + 1 + \frac{1}{x-1} \quad \Rightarrow}$\\\\
    $ \displaystyle{\quad \Rightarrow \quad x + 1 + \frac{1}{x-1} \quad \Rightarrow \quad \lim\limits_{x \to \pm \infty} \frac{1}{x-1} = 0 \quad \Rightarrow \quad f(x) = x + 1 + 0}$\\\\
\end{tabular}\\
\vspace{0.2em}
$$ \text{Asymptotengerade:} \quad y = mx + q    \qquad m \text{ und } q \text{ sind gesucht} $$

\renewcommand{\arraystretch}{1.8}
\begin{tabular}{ll}
    Steigung:           & $m = \lim \limits_{x \to \infty} \frac{f(x)}{x}$ \\
    Achsenabschnitt:    & $q = \lim \limits_{x \to \infty} (f(x) - mx)$  \quad \textrightarrow\ berechnetes $m$ einsetzen!
\end{tabular}
\renewcommand{\arraystretch}{1}